\section{The Short Game}

\begin{rulebox}{General Rule}
For players whose time is limited or for those who would like a less involved play situation, a shorter version of the game is provided. The short game begins with game turn 8 and continues through the end of game turn 12, lasting five game turns. Instead of portraying the battle situation from 4 o'clock in the afternoon (which was historically the point at which both armies finished occupying the field), the short game shows the situation as it was at 7:30 in the evening, the point at which the two sides stopped skirmishing and began to seriously fight the battle. Except as noted in this section, all the rules of 1.0 through 15.0 are in effect throughout the short game.
\end{rulebox}

\caseheading{16.1} All artillery pieces, units, and leaders (except for the Royalist leaders Rupert, Newcastle, and Eythin) are set up at the start of the game in the set-up hexes printed on their pieces. Rupert, Newcastle, and Eythin are set up in hex 2311.

\caseheading{16.2} All Royalist units with a morale rating less than 3 are set up with their back (disrupted) sides showing.

As soon as the players have set up their pieces (and before play begins), the Royalist player makes a rally check for each and every one of his disrupted units, just as if his rally phase were currently in progress. Units that pass their rally check are flipped to their front (ordered) side. Units that fail their rally check start the game disrupted.

\caseheading{16.3} The Allied Army's demoralization level is 60, and the Royalist Army's demoralization level is 50.

\caseheading{16.4} The special visibility effects (5.4) are in effect throughout game turn 8.

There is no visibility check on game turn 8. Instead, the players are to assume that an unmodified 6 was rolled during a hypothetical visibility check during step 1 of that game turn. The effects of case 5.4 never occur on any other game turn during the short game.

\caseheading{16.5} Ditch hexsides have no effect on play during game turn 8.

Ditch hexsides are ignored for both movement and combat during game turn 8. Thereafter, the ditch hexsides are treated as shown on the Terrain Effects Chart for the rest of the game. This rule represents the effect of the Allied Army's surprise charge combined with the flash rainstorm that soaked the powder of the Royalist musketeers guarding the ditch. The Allied player will tend to benefit far more than the Royalist player from this rule.
